\subsection{傅里叶变换}
\label{subsec:fourier_transform}
前面我们讨论了周期性体系的傅里叶级数展开.现在我们希望研究非周期性函数的傅里叶展开问题.
假设$f(x)$是定义在$-\infty < x < \infty$的实函数,我们依旧对该函数进行傅里叶级数展开,
该级数必须理解为周期$2\ell$扩展至$\infty$的极限情况.
由于三角函数族的变量为
\[
 \frac{n\pi x}{\ell}    
\]
引入非连续变量
\[\omega_n = \frac{n\pi} {\ell}, \quad k = 0,1,2,\cdots
    \]
可知$\omega_n = n \omega_0$, 其中$\omega_0 = \frac{\pi}{\ell}$为基础频率(Fundamental Frequency),
\[\Delta \omega_n = \omega_n - \omega_{n-1} = \omega_0 . \] 
有
\begin{equation}
    f(x) = \lim_{\ell\to \infty} \sum_{n=-\infty}^{\infty} \left[
        c_n e^{\imath \omega_n x}
        % \frac{a_0}{2} + \sum_{n=1}^{\infty} \left( a_n \cos \omega_n x + b_n \sin \omega_n x \right)
    \right]
\end{equation}
其中系数
\begin{equation}
    c_n = \frac{1}{2\ell} \int_{-\ell}^{\ell} f(x) e^{-\imath \omega_n x} dx 
\end{equation}
取$\ell\to \infty$的极限,则
\[
  \sum_n \to  \frac{ \ell}{\pi}  \int d\omega 
\]
傅里叶系数记为$F(\omega)$,
有
\begin{equation}
    f(x) = \int_{-\infty}^{\infty} d \omega F(\omega) e^{\imath \omega x}, % \quad F (\omega) \equiv 2\ell c_n
\end{equation}
上式的积分称为\textbf{傅里叶积分}(Fourier integral),是$f(x)$的\textbf{傅里叶变换}(Fourier transform).
其中
\begin{equation}
    F (\omega) \equiv \lim_{\ell\to \infty} \frac{ 2\ell c_n}{2\pi}  = \frac{1}{2\pi} \int_{-\infty}^{\infty} e^{-\imath \omega x} f(x) dx,
\end{equation}
这被称为\textbf{逆傅里叶变换}(inverse Fourier transform). 这里积分前的系数并不要紧,更关键的是积分关于$\omega$的形式.

傅里叶积分存在的条件由\textbf{傅里叶积分定理}判定: 
若函数 $f(x)$ 在区间 $(-\infty, \infty)$ 上满足条件: $(1) f(x)$ 在 任一有限区间上满足狄里希利条件; (2) $f(x)$ 在 $(-\infty, \infty)$ 
上绝对可积 (即 $\int_{-\infty}^{\infty}|f(x)| dx$ 收敛), 则 $f(x)$ 可表示成傅里叶积分, 且
$$
\textrm{傅里叶积分值} = \lim_{\epsilon\to 0} \half \left[ f(x+\epsilon) + f(x-\epsilon) . \right]
$$
% $F (\omega) \equiv 2\ell c_n$.

% \begin{align}
%     a_n =  \lim_{\ell\to \infty} \frac{1}{\ell} \int_{-\ell}^{\ell} f(x) \cos {  \omega_n x  } dx
%     \\
%     b_n = \frac{1}{\ell} \int_{-\ell}^{\ell} f(x) \sin {   \omega_n x } dx
% \end{align}

正余弦形式利用 $c_n$ 与 $a_n$、$b_n$ 的关系得到.
上一节已有 ($n>0$)
\[
  c_n=\half(a_n-\imath b_n),
  \qquad
  c_{-n}=\half(a_n+\imath b_n),
\]
亦即 $a_n=c_n+c_{-n}$, $b_n=\imath(c_n-c_{-n})$.
把 $\pm n$ 两项合并:
\[
  c_n e^{\imath\omega_n x}+c_{-n}e^{-\imath\omega_n x}
  =a_n\cos(\omega_n x)+b_n\sin(\omega_n x).
\]
又 $F(\omega)=\lim(2\ell c_n)/(2\pi)=\lim c_n/\Delta\omega$, 相应定义
\[
  A(\omega)=\lim\frac{a_n}{\Delta\omega}=F(\omega)+F(-\omega),
  \quad
  B(\omega)=\lim\frac{b_n}{\Delta\omega}=\imath\bigl[F(\omega)-F(-\omega)\bigr],
\]
或等价地 $F(\omega)=\half\bigl[A(\omega)-\imath B(\omega)\bigr]$.
由 $a_n=(1/\ell)\int f\cos$ 与 $1/\ell=\Delta\omega/\pi$, 取 $\ell\to\infty$ 得
\begin{equation}
    \left\{\begin{array}{l}
    A(\omega)=\frac{1}{\pi} \int_{-\infty}^{\infty} f(t) \cos \omega t\, d t, \\
    B(\omega)=\frac{1}{\pi} \int_{-\infty}^{\infty} f(t) \sin \omega t\, d t .
    \end{array}\right.
\end{equation}
级数中 $\pm n$ 合并后再取极限, 得到正余弦形式
\begin{equation}
    f(x) =\int_{0}^{\infty} A(\omega) \cos {\omega x} d\omega + \int_{0}^{\infty} B(\omega) \sin {\omega x} d\omega.
\end{equation}

上式还可以写成

\begin{equation}
    f(x)=\int_0^{\infty} C(\omega) \cos [\omega x-\varphi(\omega)] d \omega,
\end{equation}

其中
$$
\begin{gathered}
C(\omega)=\sqrt{[A(\omega)]^2+[B(\omega)]^2} \\
\varphi(\omega)=\arctan [B(\omega) / A(\omega)] .
\end{gathered}
$$
$C(\omega)$ 称为 $f(x)$ 的振幅谱, $\varphi(\omega)$ 称为 $f(x)$ 的相位谱.如果$f(x)$是奇函数或偶函数,
对应的级数为\textbf{傅里叶余弦积分}或\textbf{傅里叶正弦积分}.

复数形式的傅里叶积分可以写成对称的形式,
\begin{equation}
    \begin{aligned}
    f(x) & =\frac{1}{\sqrt{2 \pi}} \int_{-\infty}^{\infty} F(\omega)e^{\imath \omega x} d \omega, \\
    F(\omega) & =\frac{1}{\sqrt{2 \pi}} \int_{-\infty}^{\infty} f(x)\left[e^{\imath \omega x}\right]^* d x .
    \end{aligned}
\end{equation}
并常用符号简写为
$$
F(\omega)=\mathcal{F}[f(x)], \quad f(x)=\mathcal{F}^{-1}[F(\omega)] .
$$
$f(x)$ 和 $F(\omega)$ 分别称为傅里叶变换的\textbf{原函数}和\textbf{像函数}.

\begin{example}
求以下函数的复数傅里叶变换.
\begin{enumerate}
    \item $f(x) = e^{-\alpha |x|}, \alpha > 0$;
    \item $f(x) = \delta (x)$;
    \item $f(x) = e^{-\alpha x^2}, \alpha > 0$.
\end{enumerate}
\end{example}
\begin{solution}
\begin{enumerate}
    \item 对$x$分段进行处理,利用对称形式的傅里叶变换有
    \begin{align}
        F(\omega) &=
        \sqrt{\frac{1}{2 \pi}} \int_{-\infty}^0 e^{\alpha x+i \omega t} dx \nonumber
        % \\ &
        +
        \sqrt{\frac{1}{2 \pi}} \int_0^{\infty} e^{-\alpha x+i \omega t} dx  \nonumber
        \\
        &= \sqrt{\frac{1}{2 \pi}} \left[\frac{1}{\alpha + \imath \omega } + \frac{1}{\alpha - \imath \omega } \right] \nonumber
        \\
        &= \sqrt{\frac{1}{2 \pi}} \frac{2\alpha}{\alpha^2 + \omega^2}, \nonumber
    \end{align}
    越大的$\alpha$则$f(x)$越窄,集中在$x=0$附近,也就是越局域,而其傅里叶变换则越弥散.当$f(x)$是偶函数时,傅里叶变换的像函数为实函数.

    \item 代入
    \begin{align}
        F(\omega) &=
        \sqrt{\frac{1}{2 \pi}} \int_{-\infty}^{\infty} \delta(x) e^{+\imath \omega x} d x = \sqrt{\frac{1}{2 \pi}} ,
    \end{align}
    $f(x)$无限局域对应的傅里叶变换为无限弥散.

    \item 类似的,  
    \begin{align}
        F(\omega) &=\frac{1}{\sqrt{2 \pi}} \int_{-\infty}^{\infty} e^{-\alpha x^2} e^{\imath \omega x} d x \nonumber
        % \\ &
        % \\
        % & 
        = \frac{ e^{-\omega^2 / 4\alpha }}{\sqrt{2 \pi}} \int_{-\infty}^{\infty} e^{-\alpha  \left( x - \frac{\imath \omega}{2\alpha} \right)^2}  d x
        \nonumber
        % \\ & 
        \\
        & = \frac{ e^{-\omega^2 / 4\alpha }}{\sqrt{2 \pi}} \int_{-\infty - \imath \omega / 2\alpha}^{\infty - \imath \omega / 2\alpha} e^{-\alpha t^2} dt
        \nonumber
        % \\ &
        \\
        & = \frac{ e^{-\omega^2 / 4\alpha }}{\sqrt{2 \pi}} \sqrt{\frac{\pi}{\alpha}} 
        \nonumber
        % \\ &
        \\
            & = \frac{1}{\sqrt{2\alpha}} e^{- \frac{\omega^2}{4\alpha}} \nonumber .
        % \\ &
    \end{align}
    上式用了变量代换$t = x - \imath \omega /2 \alpha$,高斯函数的像函数还是高斯函数.
\end{enumerate}
\end{solution}


\subsection{傅里叶变换基本性质}
\idx{卷积}{juanji}
傅里叶变换满足以下基本性质.假定$f(x)$的傅里叶变换存在,记为$\mathcal{F}[f(x)] = F(\omega)$.

\begin{enumerate}
    \item 导数定理
    \begin{equation}
        \mathcal{F} [f'(x)] = \imath \omega F(\omega)
    \end{equation}

    \item 积分定理
    \begin{equation}
        \mathcal{F} [ \int^{x} f(x) dx ] = \frac{1}{\imath \omega} F(\omega)
    \end{equation}

    \item 相似性定理
    \begin{equation}
        \mathcal{F} [ f(ax) ] = \frac{1}{|a|} F\Bigl(\frac{\omega}{a}\Bigr),
        \qquad a\neq 0.
    \end{equation}
    绝对值来自雅可比. $a<0$ 时坐标反向,写成 $1/a$ 会丢一个负号.
    \item 延迟定理
    \begin{equation}
        \mathcal{F} [ f(x - x_0 ) ] = e^{-\imath \omega x_0} F(\omega),
    \end{equation}
    \item 位移定理
    \begin{equation}
        \mathcal{F} [ e^{\imath \omega_0 x} f(x) ] = F(\omega - \omega_0),
    \end{equation}
    \item 卷积定理. 卷积定义为
    \begin{equation}
        (f_1 * f_2)(x) = \int_{-\infty}^{\infty} f_1(\alpha)\, f_2(x-\alpha)\, d\alpha .
    \end{equation}
    注意第二个因子的自变量是 $x-\alpha$,不是 $\alpha-x$.在本节的对称约定下,
    \begin{equation}
        \mathcal{F}[f_1 * f_2] = \sqrt{2\pi}\, F_1(\omega)\, F_2(\omega).
    \end{equation}
    因子 $\sqrt{2\pi}$ 来自归一化.若把 $2\pi$ 全部放在逆变换上,这个常数会换成别的数,乘积结构不变.
\end{enumerate}

导数定理可以由分部积分直接看出.边界项在无穷远处消失时,
\begin{equation}
    \int_{-\infty}^{\infty} f'(x)\, e^{-\imath \omega x}\,\mathrm{d}x
    = \imath\omega \int_{-\infty}^{\infty} f(x)\, e^{-\imath \omega x}\,\mathrm{d}x,
\end{equation}
两边再除以 $\sqrt{2\pi}$ 就是 $\mathcal{F}[f']=\imath\omega F$.延迟定理是换元 $u=x-x_0$:指数多出一个 $e^{-\imath\omega x_0}$,积分变回 $F(\omega)$.

\begin{example}
求宽度为 $2a$ 的一维狭缝 $f(x)=1$（$|x|\le a$）的傅里叶变换,并由此说明矩孔.
\end{example}
\begin{solution}
\begin{equation}
    F(\omega)
    = \frac{1}{\sqrt{2\pi}}\int_{-a}^{a} e^{-\imath\omega x}\,\mathrm{d}x
    = \sqrt{\frac{2}{\pi}}\, a\, \frac{\sin(a\omega)}{a\omega}.
\end{equation}
$\omega=0$ 时商式的极限是 $\sqrt{2/\pi}\,a$.缝越宽,中央峰越窄.二维矩孔是两个这样的因子相乘.
\end{solution}

\subsection{高维傅里叶变换}

二维连续函数 $f(x, y)$ 的傅里叶变换定义如下:
设 $f(x, y)$ 是两个独立变量 $x, y$ 的函数, 且在 $\pm \infty$ 上绝对可积, 则定义积分
$$
F\left(k_1, k_2\right)=\frac{1}{2 \pi} \int_{-\infty}^{\infty} \int_{-\infty}^{\infty} f(x, y) e^{-\imath\left(k_1 x+k_2 y\right)} d x d y
$$
为二维连续函数 $f(x, y)$ 的傅里叶变换,并定义
$$
f(x, y)=\int_{-\infty}^{\infty} \int_{-\infty}^{\infty} F\left(k_1, k_2\right) e^{\imath\left(k_1 x+k_2 y\right)} d k_1 d k_2
$$
为 $F\left(k_1, k_2\right)$ 的逆变换.
这是把 $2\pi$ 全部放在正变换上的写法.与一维例子相同的对称约定是两侧各有一个 $1/(2\pi)$:
\begin{equation}
    F(k_1,k_2)=\frac{1}{2\pi}\iint f(x,y)\,e^{-\imath(k_1 x+k_2 y)}\,\mathrm{d}x\,\mathrm{d}y,
\end{equation}
\begin{equation}
    f(x,y)=\frac{1}{2\pi}\iint F(k_1,k_2)\,e^{\imath(k_1 x+k_2 y)}\,\mathrm{d}k_1\,\mathrm{d}k_2.
\end{equation}
两种约定不能在同一次计算里混用.下面的矩孔按对称约定计算.

\begin{example}
    求函数 $f(x, y)= \begin{cases}A, & |x| \leq X,|y| \leq Y \\ 0, & |x|>X,|y|>Y\end{cases}$ 的傅里叶变换 (矩孔费琅和夫衍射).
\end{example}
\begin{solution}
    由傅里叶变换关系
$$
F\left(k_1, k_2\right)=\frac{1}{2 \pi} \int_{-\infty}^{\infty} \int_{-\infty}^{\infty} f(x, y) e^{-i\left(k_1 x+k_2 y\right)} d x d y
$$
有
$$
\begin{aligned}
F\left(k_1, k_2\right) & =\frac{A}{2 \pi} \int_{-X}^X e^{-i k_1 x} d x \int_{-Y}^Y e^{-i k_2 y} d y \\
& =\left.\left.\frac{A}{2 \pi} \frac{1}{-i k_1} e^{-i k_1 x}\right|_{-X} ^X \frac{1}{-i k_2} e^{-i k_2 y}\right|_{-Y} ^Y \\
& =\frac{2A}{\pi} \frac{1}{i 2 k_1}\left(e^{i k_1 X}-e^{-i k_1 X}\right) \frac{1}{i 2 k_2}\left(e^{i k_2 Y}-e^{-i k_2 Y}\right) \\
& =\frac{2A X Y}{\pi} \frac{\sin \left(k_1 X\right)}{k_1 X} \frac{\sin \left(k_2 Y\right)}{k_2 Y}
\end{aligned}
$$
\end{solution}

物理上,$f(x,y)$ 是孔径面上的复振幅透过率:矩形开口内均匀照明取常值 $A$,挡光处为零.
费琅和夫(远场)衍射振幅正是孔径函数的傅里叶变换 $F(k_1,k_2)$,观察屏上的光强正比于 $|F|^2$;
$k_1,k_2$ 对应衍射角方向.
轴上 $F(0,0)=2AXY/\pi$,与孔径面积成正比.
$k_1 X=n\pi$ 或 $k_2 Y=m\pi$($n,m=\pm 1,\pm 2,\ldots$) 处振幅为零,屏上出现互相垂直的暗纹.
孔越窄,中央峰在对应方向越宽,与一维狭缝同一条宽窄互换.
$F$ 又是两个一维狭缝因子的乘积,故矩孔图样是互相垂直的两条单缝衍射结构.

三维时对称约定是两侧各有 $(2\pi)^{-3/2}$:
\begin{equation}
    F(\vec{k})=\frac{1}{(2 \pi)^{3/2}} \iiint_{-\infty}^{\infty} f(\vec{r}) e^{-\imath  \vec{k} \cdot \vec{r}} d^3 \vec{r},
\end{equation}
\begin{equation}
     f(\vec{r})=\frac{1}{(2 \pi)^{3/2}}\iiint_{-\infty}^{\infty} F(\vec{k}) e^{\imath \vec{k} \cdot \vec{r}} d^3 \vec{k}.
\end{equation}
把 $(2\pi)^{-3}$ 只放在一侧是另一种等价约定,不能和上式混用.指数可分离时,积分拆成三个一维积分.

若 $f(\vec{r})=f(r)$ 球对称,则 $F$ 只依赖 $k=|\vec{k}|$.
取极轴沿 $\vec{k}$,角积分给出 $4\pi\sin(kr)/(kr)$,对称约定下
\begin{equation}
    F(k)
    =\sqrt{\frac{2}{\pi}}\,\frac{1}{k}
      \int_0^{\infty} r\,f(r)\sin(kr)\,\mathrm{d}r.
\end{equation}
三维球对称傅里叶变换化成一维正弦变换.

\begin{example}
求汤川势 $f(r)=e^{-\mu r}/r$($\mu>0$) 的三维傅里叶变换,并说明它在散射中的用法.
\end{example}
\begin{solution}
用 $\int_0^{\infty}e^{-\mu r}\sin(kr)\,\mathrm{d}r=k/(k^2+\mu^2)$,得
\begin{equation}
    F(k)=\sqrt{\frac{2}{\pi}}\,\frac{1}{k^2+\mu^2}.
\end{equation}
散射计算里动量空间矩阵元常不带 $(2\pi)^{-3/2}$,直接写
\begin{equation}
    \int\frac{e^{-\mu r}}{r}\,e^{-\imath\vec{q}\cdot\vec{r}}\,\mathrm{d}^3\vec{r}
    =\frac{4\pi}{q^2+\mu^2},
    \qquad \vec{q}=\vec{k}'-\vec{k}.
\end{equation}
这是玻恩近似下一阶散射振幅的势因子.$\mu\to 0$ 给出库仑势的 $4\pi/q^2$,对应卢瑟福公式的动量依赖.
\end{solution}